\section{Introduction}

Contrastive image-text encoders such as CLIP have become default infrastructure for retrieval, zero-shot classification, and as frozen components inside larger multimodal systems~\cite{radford2021clip,ilharco2021openclip}. Their strength is broad alignment: a single embedding space makes many downstream comparisons cheap. The same coarse alignment can be a weakness for relational edits. Captions that preserve object vocabulary but swap roles, attributes, spatial relations, or word order may remain close in embedding space even when the correct image-text match should change~\cite{yuksekgonul2022bags,hsieh2023sugarcrepe,thrush2022winoground}.

This paper asks a narrow question: can a frozen CLIP encoder be made more sensitive to relational counterfactuals while keeping the pretrained embedding path dominant? We use ``guardrail'' only in this narrow technical sense: retention of retrieval and zero-shot behavior on held-out COCO and ImageNet checks. It does not refer to content safety, red-teaming, or policy filtering. Because the COCO check uses a 1,000-image subset, it serves as a signal for catastrophic retrieval collapse rather than a certificate of full retrieval parity.

The word ``dominant'' is central to the design. If a method obtains relation sensitivity by replacing the frozen embedding with a new representation, then a high relation score may only show that the adaptation has learned the staged benchmark. It does not show that the model still behaves like a useful CLIP encoder. In this project, the intended intervention is smaller: keep the original OpenCLIP representation available to every downstream dot product, and ask whether a relation-trained correction can move difficult counterfactual pairs just enough to change their ranking. This makes the evaluation asymmetric. A relation gain is necessary, but it is not sufficient. The gain must be read together with retrieval, zero-shot, residual-size, and external-validity checks.

This emphasis also determines how we interpret negative results. A hard bottleneck that improves relation accuracy but breaks COCO retrieval is not a partial success for the paper's claim; it is evidence that the relation signal is easy to overfit and difficult to add safely. A low residual norm is not automatically safe either, because two residuals with similar norm can point in different directions in embedding space and have different retrieval effects. The paper therefore treats each run as a point in a relation-preservation tradeoff rather than as a single scalar score.

The difficulty is visible in the project's early variants. Hard token bottlenecks can expose relation signal, but replacing the pretrained path damages broad CLIP behavior. A representative hard-token targeted run reaches relation accuracy $0.7488$, yet falls to COCO image-to-text R@1 $0.3920$ and ImageNet top-1 $0.0660$. In contrast, the selected RRB checkpoint reaches relation accuracy $0.7734$ with COCO R@1 $0.7340$ and ImageNet top-1 $0.6342$. The result is not lossless preservation, but it is a qualitatively different tradeoff.

The same caution applies to mechanism language. The implementation uses multiple latent residual channels, but the responsibility probes later in the project do not justify a claim that those channels become human-interpretable relation slots. We therefore use the channels as an architectural bottleneck: they constrain how the residual is computed from frozen token features, but the paper does not ask readers to believe that channel one means spatial relations or channel two means attributes. This distinction matters because the no-channel controls are relation-strong yet unsafe, while the responsibility probes are diffuse. The supported conclusion is about the residual interface and its guardrail behavior, not about discovered symbolic structure.

The paper is therefore a bottleneck-first study rather than a claim about rank-efficient adaptation in general. LoRA-style methods constrain parameter updates through rank, while RRB constrains the route from frozen token features to the final embedding correction. This difference is the object of study: we evaluate whether a narrow residual interface can carry relation supervision without letting the adapted path dominate the pretrained embedding.

We make four contributions.
\begin{enumerate}
    \item We define Relational Residual Bottlenecking (RRB), a parser-free residual adaptation module for frozen OpenCLIP image and text encoders.
    \item We show that RRB improves staged relation accuracy while avoiding the catastrophic retrieval and zero-shot collapse of hard bottlenecks and no-channel residual controls.
    \item We isolate the useful training signal: structured relational counterfactuals under residual anchoring, not generic hard negatives and not the earlier targeted relational loss.
    \item We report the limits directly: COCO retention is seed-fragile, Winoground is not improved, Flickr30k image-to-text retrieval drops, and responsibility probes do not justify semantic slot claims.
\end{enumerate}

\begin{figure*}[t]
    \centering
    % Auto-generated by figures/generate_paper_assets.py
\begin{tikzpicture}[font=\footnotesize]
\draw[->] (0,0) -- (7.05,0) node[right] {Relation accuracy};
\draw[->] (0,0) -- (0,4.25) node[above] {COCO I2T R@1};
\draw (0.40,0.04) -- (0.40,-0.04) node[below] {0.65};
\draw (2.40,0.04) -- (2.40,-0.04) node[below] {0.70};
\draw (4.40,0.04) -- (4.40,-0.04) node[below] {0.75};
\draw (6.40,0.04) -- (6.40,-0.04) node[below] {0.80};
\draw (0.04,0.50) -- (-0.04,0.50) node[left] {0.40};
\draw (0.04,1.33) -- (-0.04,1.33) node[left] {0.50};
\draw (0.04,2.17) -- (-0.04,2.17) node[left] {0.60};
\draw (0.04,3.00) -- (-0.04,3.00) node[left] {0.70};
\draw (0.04,3.83) -- (-0.04,3.83) node[left] {0.80};
\draw[densely dashed,gray] (2.22,0) -- (2.22,4.00);
\draw[densely dashed,gray] (0,3.00) -- (6.80,3.00);
\node[gray,anchor=south west] at (2.27,0.05) {R005 relation};
\node[gray,anchor=south east] at (6.75,3.05) {COCO gate};
\node[circle,draw=black,fill=black,inner sep=1.8pt] (R003) at (0.66,3.83) {};
\node[anchor=west] at (0.74,3.93) {R003};
\node[circle,draw=gray,fill=gray,inner sep=1.8pt] (R005) at (2.22,3.66) {};
\node[anchor=west] at (2.30,3.76) {R005};
\node[rectangle,draw=red!75!black,fill=red!75!black,inner sep=1.8pt] (R011) at (4.35,0.43) {};
\node[anchor=west] at (4.43,0.53) {R011};
\node[rectangle,draw=red!75!black,fill=red!75!black,inner sep=1.8pt] (R014) at (3.73,0.59) {};
\node[circle,draw=blue!75!black,fill=blue!75!black,inner sep=1.8pt] (R019) at (5.19,3.16) {};
\node[circle,draw=blue!75!black,fill=blue!75!black,inner sep=1.8pt] (R023) at (6.09,2.85) {};
\node[circle,draw=gray,fill=gray,inner sep=1.8pt] (R024) at (1.34,3.83) {};
\node[circle,draw=blue!75!black,fill=blue!75!black,inner sep=2.3pt] (R028) at (5.34,3.28) {};
\node[anchor=east] at (5.26,3.38) {R028};
\node[circle,draw=blue!75!black,fill=blue!75!black,inner sep=1.8pt] (R029) at (5.32,3.23) {};
\node[circle,draw=blue!75!black,fill=blue!75!black,inner sep=1.8pt] (R031) at (5.64,3.22) {};
\node[circle,draw=blue!75!black,fill=blue!75!black,inner sep=1.8pt] (R032) at (5.62,2.78) {};
\node[diamond,draw=orange!85!black,fill=orange!85!black,inner sep=1.8pt] (R033) at (5.70,0.28) {};
\node[anchor=east] at (5.62,0.38) {R033};
\node[circle,draw=cyan!60!black,fill=cyan!60!black,inner sep=1.8pt] (R035) at (5.61,3.11) {};
\node[circle,draw=cyan!60!black,fill=cyan!60!black,inner sep=1.8pt] (R036) at (5.60,2.85) {};
\node[circle,draw=cyan!60!black,fill=cyan!60!black,inner sep=1.8pt] (R037) at (5.66,3.12) {};
\node[circle,draw=cyan!60!black,fill=cyan!60!black,inner sep=1.8pt] (R038) at (5.34,3.36) {};
\node[anchor=east] at (5.26,3.46) {R038};
\node[circle,draw=cyan!60!black,fill=cyan!60!black,inner sep=1.8pt] (R039) at (5.58,2.95) {};
\node[diamond,draw=orange!85!black,fill=orange!85!black,inner sep=1.8pt] (R040) at (5.49,1.18) {};
\node[circle,draw=purple!75!black,fill=purple!75!black,inner sep=1.8pt] (R043) at (4.17,3.30) {};
\node[anchor=west] at (4.25,3.40) {R043};
\draw[dashed,thick,black!55] (1.34,3.83) -- (2.22,3.66) -- (5.34,3.36) -- (6.09,2.85);
\node[black!60,anchor=west] at (4.65,2.70) {evaluated frontier};
\node[anchor=west] at (0.15,4.05) {\tikz\node[circle,draw=blue!75!black,fill=blue!75!black,inner sep=1.5pt]{}; RRB};
\node[anchor=west] at (1.55,4.05) {\tikz\node[rectangle,draw=red!75!black,fill=red!75!black,inner sep=1.4pt]{}; hard token};
\node[anchor=west] at (3.25,4.05) {\tikz\node[diamond,draw=orange!85!black,fill=orange!85!black,inner sep=1.5pt]{}; no-channel};
\node[anchor=west] at (5.25,4.05) {\tikz\node[circle,draw=cyan!60!black,fill=cyan!60!black,inner sep=1.5pt]{}; anchor-20};
\end{tikzpicture}

    \caption{Relation accuracy versus COCO image-to-text retrieval for key evaluated systems. Hard-token controls replace the pretrained embedding path with a bottlenecked representation, while no-channel controls use a residual path without the multi-channel bottleneck. The selected RRB checkpoint is shown as a point; seed dispersion for RRB variants is summarized separately in Figure~\ref{fig:seed_anchor}. Standard PEFT baselines such as LoRA and text-only tuning were not evaluated in this ablation study, which focuses on the architectural impact of the RRB bottleneck.}
    \Description{Scatter plot of relation accuracy against COCO image-to-text recall. OpenCLIP and the plain adapter sit at lower relation and high retrieval, hard-token and no-channel controls have higher relation but poor retrieval, and the selected RRB point sits between them with improved relation and bounded retrieval loss.}
    \label{fig:relation_coco}
\end{figure*}

This framing is intentionally conservative. The paper is not a claim that RRB solves compositional reasoning for CLIP. It is an empirical characterization of a useful relation-vs-preservation tradeoff and of the failure modes that remain.

The rest of the paper follows that framing. Section~2 positions RRB relative to compositional CLIP evaluation, structure-aware retrieval, and preservation-oriented adaptation. Section~3 defines the residual interface and training objective. Section~4 reports the run families and shows why the selected checkpoint is a tradeoff point rather than a universal improvement. Section~5 analyzes the remaining failures, including seed fragility, external retrieval degradation, and unsupported channel-interpretability claims.
